% Landscape Poster: an A0 landscape poster on a four-column grid.
% Every block is placed on the grid with textpos, so blocks never push each other.
% Edit the metadata and colours below, then the blocks. Build: latexmk -pdf main.tex
\documentclass{article}
\usepackage[paperwidth=1189mm,paperheight=841mm,margin=0mm]{geometry}
\usepackage[T1]{fontenc}
% Roboto for text, Roboto Slab (\robotoslab) for headings; lining tabular figures
% keep numbers aligned in the table and the charts.
\usepackage[sfdefault,lining,tabular]{roboto}
\usepackage{newtxsf}
\usepackage[letterspace=120]{microtype}
\usepackage{booktabs,tabularx,multicol,qrcode,xcolor}
\usepackage[absolute,overlay]{textpos}
\usepackage[hidelinks]{hyperref}
\usepackage{tikz,pgfplots}
\usetikzlibrary{calc,arrows.meta,positioning}
\pgfplotsset{compat=1.18}
\usepackage[style=numeric-comp,sorting=none,maxbibnames=1,minbibnames=1,
  giveninits=true,doi=false,url=false,isbn=false,backend=biber]{biblatex}
\addbibresource{refs.bib}
% Compact references: no titles of journals in italics clutter, no "In:".
\renewbibmacro{in:}{}
\AtEveryBibitem{\clearfield{number}\clearfield{month}}

% ---------------------------------------------------------------- colours
% lsCobalt carries headings, squares and the study's main colour.
% lsAmber is reserved for the one data series the poster is about.
\definecolor{lsCobalt}{HTML}{1F43A6}
\definecolor{lsAmber}{HTML}{C99234}
\definecolor{lsInk}{HTML}{1D2230}
\definecolor{lsMuted}{HTML}{687080}
\definecolor{lsRule}{HTML}{D8D3C8}
\definecolor{lsPaper}{HTML}{FBF9F4}
\colorlet{lsPanel}{lsCobalt!5!lsPaper}
\colorlet{lsGrey}{lsMuted!55!lsPaper}
\pagecolor{lsPaper}\color{lsInk}

% ---------------------------------------------------------------- metadata
\newcommand{\postertitle}{Later bells, longer sleep}
\newcommand{\postersubtitle}{A cluster-randomised trial of school start times in 24 secondary schools}
\newcommand{\posterauthors}{Ines Varga\textsuperscript{1}\quad Daniel Osei\textsuperscript{1,2}\\Hanna Lindqvist\textsuperscript{2}\quad Marco Bellini\textsuperscript{1}}
\newcommand{\posteraffiliations}{\textsuperscript{1}\,Centre for Sleep and Learning, Eastmoor University\par
  \textsuperscript{2}\,Department of Paediatrics, Halden University Hospital}
\newcommand{\postervenue}{European Sleep Research Congress 2026 \textperiodcentered{} Poster B-27}
\newcommand{\posterurl}{https://example.org/later-bells}
\newcommand{\postercontact}{i.varga@example.org}

% ---------------------------------------------------------------- type scale
% Sizes for reading at 1.5 m; leading is about 1.4 times the size.
\renewcommand{\normalsize}{\fontsize{34}{47}\selectfont}
\renewcommand{\small}{\fontsize{28}{38}\selectfont}
\renewcommand{\footnotesize}{\fontsize{23}{31}\selectfont}
\newcommand{\lsmedium}{\fontseries{medium}\selectfont}
\newcommand{\lsslab}[1]{{\robotoslab\bfseries#1}}
\newcommand{\lskicker}[1]{{\footnotesize\lsmedium\textls[160]{\MakeUppercase{#1}}}}
\setlength{\parindent}{0pt}
\setlength{\parskip}{0.5em}
\pagestyle{empty}
\renewcommand*{\bibfont}{\fontsize{20}{26}\selectfont}
\setlength{\bibitemsep}{0.3em}
\setlength{\biblabelsep}{4mm}

% ---------------------------------------------------------------- grid
% 40 mm margins, four columns, 30 mm gutters. \lsx{k} is the left edge of column k.
\newlength{\lsmargin}\setlength{\lsmargin}{40mm}
\newlength{\lsgutter}\setlength{\lsgutter}{30mm}
\newlength{\lscol}\setlength{\lscol}{\dimexpr(\paperwidth-2\lsmargin-3\lsgutter)/4\relax}
\newlength{\lswide}\setlength{\lswide}{\dimexpr2\lscol+\lsgutter\relax}
\newcommand{\lsx}[1]{\dimexpr\lsmargin+#1\lscol+#1\lsgutter-\lscol-\lsgutter\relax}
\setlength{\TPHorizModule}{1mm}\setlength{\TPVertModule}{1mm}
\newlength{\lsbody}\setlength{\lsbody}{176mm}   % top of the body row
\newlength{\lsfoot}\setlength{\lsfoot}{722mm}   % top of the footer row

% ---------------------------------------------------------------- headings
% A section opens with a small cobalt square holding its number.
\newcommand{\lssection}[2]{%
  \par\addvspace{14mm}%
  \begin{tikzpicture}[baseline=(t.base)]
    \node[fill=lsCobalt,text=lsPaper,minimum size=18mm,inner sep=0,
      font=\fontsize{29}{29}\selectfont\bfseries] (n) {#1};
    \node[anchor=west,inner sep=0,text=lsInk,font=\fontsize{60}{66}\selectfont]
      (t) at ([xshift=8mm]n.east) {\lsslab{#2}};
  \end{tikzpicture}\par\vspace{4mm}}
\newcommand{\lssub}[1]{\par\vspace{2mm}{\lsmedium#1\par}\vspace{1mm}}
\newcommand{\lscaption}[1]{\par\vspace{2mm}{\footnotesize\color{lsMuted}#1\par}}

\pgfplotsset{lsaxis/.style={
  axis lines=left, axis line style={lsMuted,line width=1pt},
  tick style={lsMuted,line width=1pt}, major tick length=2.5mm,
  tick label style={font=\footnotesize,text=lsMuted},
  label style={font=\footnotesize,text=lsMuted},
  grid=major, major grid style={lsRule,line width=0.7pt},
  scale only axis, clip=false}}

\begin{document}
% ================================================================ header
\begin{tikzpicture}[remember picture,overlay]
  \fill[lsCobalt] (current page.north west) rectangle ([yshift=-9mm]current page.north east);
  \draw[lsRule,line width=1.2pt] ([shift={(\lsmargin,-150mm)}]current page.north west)
    -- ([shift={(-\lsmargin,-150mm)}]current page.north east);
\end{tikzpicture}

\begin{textblock*}{\dimexpr\lswide+\lscol\relax}(\lsmargin,34mm)
  {\color{lsCobalt}\lskicker{\postervenue}}\par\vspace{6mm}
  {\fontsize{112}{118}\selectfont\color{lsInk}\lsslab{\postertitle}\par}\vspace{7mm}
  {\fontsize{38}{44}\selectfont\robotoslab\color{lsMuted}\postersubtitle\par}
\end{textblock*}

\begin{textblock*}{\lscol}(\lsx{4},36mm)
  \raggedright
  {\small\lsmedium\setlength{\tabcolsep}{0pt}\posterauthors\par}\vspace{3mm}
  {\footnotesize\color{lsMuted}\posteraffiliations\par}\vspace{9mm}
  % Placeholder logos: replace each picture with \includegraphics[height=24mm]{logo}.
  \begin{tikzpicture}[x=1mm,y=1mm]
    \draw[lsCobalt,line width=3pt] (15,15) circle (14);
    \fill[lsCobalt] (15,15) -- ++(90:14) arc (90:210:14) -- cycle;
    \node[anchor=west,inner sep=0,align=left,font=\fontsize{20}{23}\selectfont,text=lsInk]
      at (35,15) {\lsslab{Eastmoor}\\University};
    \begin{scope}[xshift=118mm]
      \fill[lsInk] (0,1) rectangle ++(28,28);
      \fill[lsPaper] (6,7) rectangle ++(6,15) (15,7) rectangle ++(6,6);
      \node[anchor=west,inner sep=0,align=left,font=\fontsize{20}{23}\selectfont,text=lsInk]
        at (35,15) {\lsslab{Halden}\\Hospital};
    \end{scope}
  \end{tikzpicture}
\end{textblock*}

% ================================================================ column 1
\begin{textblock*}{\lscol}(\lsx{1},\lsbody)
\raggedright
\lssection{1}{Background}
Adolescents fall asleep about two hours later than children~\cite{roenneberg2004},
yet most schools start before 08:30. The result is chronic short sleep on school
nights~\cite{carskadon2011}.

Observational studies link later starts to longer sleep~\cite{wahlstrom2002,dunster2018},
but randomised evidence is scarce~\cite{minges2016}.

\lssection{2}{Study design}
\begin{tikzpicture}[x=1mm,y=1mm,
    box/.style={draw=lsCobalt,line width=1.6pt,fill=lsPaper,inner xsep=5mm,inner ysep=9mm,
      text width=#1,align=center,font=\small},
    box/.default=242mm,
    arm/.style={box=64mm},
    flow/.style={lsCobalt,line width=1.6pt,-{Latex[length=4mm,width=4mm]}}]
  \node[box] (a) at (127.4,0) {{\lsmedium 24 schools}, 3\,412 pupils aged 13--16};
  \node[box,below=18mm of a] (b) {Baseline: 2 weeks of wrist actigraphy and sleep diaries};
  \node[box,below=18mm of b,fill=lsCobalt,text=lsPaper] (c) {Randomised by school, 1\,:\,1\,:\,1};
  \node[arm,below=30mm of c] (m) {{\lsmedium 08:30}\\8 schools};
  \node[arm,left=12mm of m] (l) {{\lsmedium 08:00}\\control};
  \node[arm,right=12mm of m,draw=lsAmber,line width=3pt] (r) {{\lsmedium 09:00}\\8 schools};
  \node[box,below=30mm of m] (d) {Week 14: actigraphy, diaries, attendance, maths test};
  \draw[flow] (a) -- (b);
  \draw[flow] (b) -- (c);
  \foreach \k in {l,m,r}{\draw[flow] (c.south) -- (\k.north);
    \draw[flow] (\k.south) -- (\k.south|-d.north);}
\end{tikzpicture}
\lscaption{Schools, not pupils, were randomised, so timetables changed for
everyone. Analysis by intention to treat with a random intercept per school.}
\end{textblock*}

% ================================================================ results panel
\begin{textblock*}{\lswide}(\lsx{2},\dimexpr\lsbody-12mm\relax)
\begin{tikzpicture}[remember picture]
  \fill[lsPanel] (0,0) rectangle (\lswide,-532mm);
  \fill[lsCobalt] (0,0) rectangle (\lswide,-2.4mm);
\end{tikzpicture}
\end{textblock*}

\begin{textblock*}{\dimexpr\lswide-36mm\relax}(\dimexpr\lsx{2}+18mm\relax,\lsbody)
\lssection{3}{Results}
\begin{minipage}[t]{\dimexpr0.5\linewidth-12mm\relax}
  \lssub{Sleep on school nights}
  \begin{tikzpicture}
    \begin{axis}[lsaxis, width=\linewidth-70mm, height=215mm,
        xmin=0, xmax=14, ymin=6.6, ymax=8.2, xtick={0,2,...,14}, ytick={6.6,7,...,8.2},
        y tick label style={/pgf/number format/.cd,fixed,fixed zerofill,precision=1},
        xlabel={Week of term}, ylabel={Mean sleep (hours)},
        every axis plot/.append style={line width=2.6pt}]
      \addplot[lsGrey] coordinates {(0,6.92)(1,6.95)(2,6.90)(3,6.94)(4,6.88)(5,6.91)(6,6.93)
        (7,6.89)(8,6.92)(9,6.90)(10,6.87)(11,6.91)(12,6.89)(13,6.90)(14,6.88)}
        node[right,font=\footnotesize,text=lsMuted] {08:00};
      \addplot[lsCobalt] coordinates {(0,6.95)(1,7.18)(2,7.25)(3,7.30)(4,7.27)(5,7.33)(6,7.31)
        (7,7.35)(8,7.32)(9,7.36)(10,7.34)(11,7.37)(12,7.35)(13,7.38)(14,7.36)}
        node[right,font=\footnotesize,text=lsCobalt] {08:30};
      \addplot[lsAmber,line width=4.4pt] coordinates {(0,6.90)(1,7.42)(2,7.58)(3,7.66)(4,7.70)
        (5,7.68)(6,7.74)(7,7.71)(8,7.76)(9,7.73)(10,7.78)(11,7.75)(12,7.79)(13,7.77)(14,7.80)}
        node[right,font=\footnotesize\bfseries,text=lsAmber!80!lsInk] {09:00};
      \draw[lsMuted,line width=1pt,dashed] (axis cs:0.5,6.6) -- (axis cs:0.5,8.2)
        node[pos=0.97,right,font=\footnotesize,text=lsMuted] {new start time};
    \end{axis}
  \end{tikzpicture}
  \lscaption{Weekly means from actigraphy, all pupils. The 09:00 arm gains
  54 minutes within three weeks and keeps them.}
\end{minipage}\hfill
\begin{minipage}[t]{\dimexpr0.5\linewidth-12mm\relax}
  \lssub{Pupils sleeping at least 8 hours}
  \begin{tikzpicture}
    \begin{axis}[lsaxis, ybar, width=\linewidth-42mm, height=215mm,
        bar width=15mm, area legend, ymin=0, ymax=60, ytick={0,20,40,60},
        symbolic x coords={Year 8,Year 9,Year 10,Year 11}, xtick=data,
        enlarge x limits=0.16, xmajorgrids=false, major x tick style={draw=none},
        x tick label style={font=\footnotesize,text=lsInk}, ylabel={Share of pupils (\%)},
        legend style={draw=none,fill=none,font=\footnotesize,at={(0,1.03)},anchor=south west,
          legend columns=3,/tikz/every odd column/.append style={column sep=2mm},
          /tikz/every even column/.append style={column sep=7mm}},
        legend image code/.code={\fill[#1] (0mm,-2mm) rectangle (8mm,2mm);}]
      \addplot[fill=lsGrey,draw=none] coordinates {(Year 8,34)(Year 9,27)(Year 10,21)(Year 11,18)};
      \addplot[fill=lsCobalt,draw=none] coordinates {(Year 8,41)(Year 9,36)(Year 10,30)(Year 11,27)};
      \addplot[fill=lsAmber,draw=none] coordinates {(Year 8,52)(Year 9,49)(Year 10,44)(Year 11,41)};
      \legend{08:00,08:30,09:00}
    \end{axis}
  \end{tikzpicture}
  \lscaption{Week 14, school nights. The gap grows with age, when the body
  clock runs latest.}
\end{minipage}

\vspace{16mm}
\begin{minipage}[t]{\dimexpr0.66\linewidth\relax}
  \lssub{Secondary outcomes at week 14, 09:00 against 08:00}
  {\small\renewcommand{\arraystretch}{1.3}%
  \begin{tabularx}{\linewidth}{@{}X r r r r@{}}
    \toprule
    {\lsmedium Outcome} & {\lsmedium 08:00} & {\lsmedium 09:00} & {\lsmedium Difference} & {\lsmedium 95\% CI} \\
    \midrule
    Sleep onset (clock time)  & 23:21 & 23:34 & $+13$ min  & $+4$ to $+22$ \\
    Wake time (clock time)    & 06:38 & 07:39 & $+61$ min  & $+52$ to $+70$ \\
    Daytime sleepiness, PDSS  & 17.8  & 14.6  & $-3.2$\phantom{ min}  & $-4.1$ to $-2.3$ \\
    Attendance (\%)           & 92.1  & 94.0  & $+1.9$\phantom{ min}  & $+0.8$ to $+3.0$ \\
    Maths test ($z$-score)    & 0.00  & 0.11  & $+0.11$\phantom{ min} & $+0.02$ to $+0.20$ \\
    \bottomrule
  \end{tabularx}\par}
  \lscaption{Mixed-effects estimates adjusted for baseline value, year group and sex;
  $n = 3\,187$ pupils with follow-up data.}
\end{minipage}\hfill
\begin{minipage}[t]{\dimexpr0.28\linewidth\relax}
  \lssub{Headline}
  {\fontsize{150}{150}\selectfont\color{lsAmber!85!lsInk}\lsslab{+54}\par}\vspace{3mm}
  {\small\raggedright minutes of sleep per school night at 09:00, against 08:00. Sleep onset
  barely moved, so the gain is almost all wake time.\par}
\end{minipage}
\end{textblock*}

% ================================================================ column 4
\begin{textblock*}{\lscol}(\lsx{4},\lsbody)
\raggedright
\lssection{4}{Take-home messages}
\begin{tikzpicture}
  \node[fill=lsCobalt,text=lsPaper,anchor=north west,inner xsep=12mm,inner ysep=14mm,
    text width=\lscol-24mm] at (0,0) {%
    \foreach \n/\msg in {
      1/{A 09:00 start gives teenagers almost an hour more sleep per night.},
      2/{The gain appears in the first three weeks and holds for the whole term.},
      3/{Pupils are less sleepy, attend more and score slightly higher in maths.}}{%
      \begin{tabularx}{\linewidth}{@{}l@{\hspace{6mm}}X@{}}
        {\fontsize{48}{48}\selectfont\lsslab{\n}} & \raggedright\arraybackslash\fontsize{37}{50}\selectfont\lsmedium\msg \\
      \end{tabularx}\par\vspace{10mm}}%
    \vspace{-10mm}};
\end{tikzpicture}

\lssection{5}{Limits and next steps}
Pupils and teachers knew the start time, so diaries may be biased; actigraphy
is not.

One term cannot show lasting effects. A two-year follow-up of all 24 schools
reports in 2028.
\end{textblock*}

% Anchored by its bottom edge so it lines up with the foot of the results panel.
\begin{textblock*}{\lscol}[0,1](\lsx{4},\dimexpr\lsbody+520mm\relax)
\raggedright
\lssection{6}{Acknowledgements}
{\small We thank the 24 schools, their pupils and families. Funded by the
Eastmoor Health Research Fund (grant EH-2291). Trial registration ISRCTN 40021873.\par}
\end{textblock*}

% ================================================================ footer
\begin{textblock*}{\dimexpr\paperwidth-2\lsmargin\relax}(\lsmargin,\lsfoot)
  {\color{lsRule}\rule{\linewidth}{1.2pt}}
\end{textblock*}

\begin{textblock*}{\dimexpr\lscol+\lswide+\lsgutter\relax}(\lsx{1},\dimexpr\lsfoot+10mm\relax)
  {\color{lsCobalt}\lskicker{References}}\par\vspace{2mm}
  \setlength{\columnsep}{\lsgutter}
  \begin{multicols}{3}\raggedright
    \printbibliography[heading=none]
  \end{multicols}
\end{textblock*}

\begin{textblock*}{\lscol}(\lsx{4},\dimexpr\lsfoot+10mm\relax)
  \begin{minipage}[t]{\dimexpr\lscol-70mm\relax}\raggedright
    {\color{lsCobalt}\lskicker{Contact}}\par\vspace{2mm}
    {\small\postercontact\par}\vspace{6mm}
    {\color{lsCobalt}\lskicker{Full paper and data}}\par\vspace{2mm}
    {\footnotesize\posterurl\par}
  \end{minipage}\hfill
  \begin{minipage}[t]{58mm}
    \vspace{0pt}{\color{lsInk}\qrcode[height=58mm]{\posterurl}}
  \end{minipage}
\end{textblock*}

\null
\end{document}
